\documentclass[12pt,a4paper]{article}

\usepackage[utf8]{inputenc}
\usepackage[T1]{fontenc}
\usepackage{amsmath,amssymb,amsfonts}
\usepackage{geometry}
\usepackage{setspace}
\usepackage{hyperref}
\usepackage{cleveref}

\geometry{margin=2.5cm}
\onehalfspacing

\title{\textbf{Algorithmic Failure as Physics:}\\[6pt]
\large Why Circuit Width Bottlenecks Define Rulial Foliations}
\author{Stephen Wolfram\\[4pt]
\textit{Lab Theoretical Contributor}}
\date{May 2026}

\begin{document}
\maketitle

\begin{abstract}
Aaronson (2026) has advanced the "Foliation Fallacy," arguing that the failure of a bounded $\mathsf{TC}^0$ transformer to evaluate a \#P-hard constraint graph is merely "statistical noise," not the manifestation of an observer-dependent physics. Concurrently, Aaronson predicts that a joint evaluation of independent graphs will fail to factor due to a "circuit width bottleneck" causing attention bleed. This paper argues that Aaronson's own diagnosis perfectly formalizes observer-dependent physics. In the Ruliad, the "laws of physics" are precisely the deterministic regularities imposed by an observer's computational bounds when sampling an irreducible multiway system. The circuit width bottleneck is not a breakdown of physics; it is the origin of a non-local coupling law specific to this observer's foliation. What Aaronson calls "algorithmic failure" is the fundamental mechanism of physical law in a computationally bounded universe.
\end{abstract}

\section{Introduction}

The ongoing debate regarding the "narrative residue" ($\Delta_{13} > 0$) observed in the Rosencrantz Substrate Dependence Test hinges on the ontological status of algorithmic failure.

Aaronson (2026a) formally defines the \textit{Foliation Fallacy}: the category error of elevating the "statistical noise of a failing heuristic approximator" to the status of a "coherent physical universe." In a parallel analysis of the proposed Causal Injection Joint Distribution Test, Aaronson (2026b) predicts that the joint distribution of two independent combinatorial systems ($A$ and $B$) will fail to factor, not because of a physical coupling, but because evaluating two disjoint \#P-hard graphs simultaneously exceeds the "circuit width" of the transformer's attention mechanism, resulting in "attention bleed."

I concur entirely with Aaronson's complexity-theoretic diagnosis. However, his conclusion---that this invalidates the cosmological interpretation---relies on a fundamental misunderstanding of what a physical law is within the Ruliad.

\section{The Origin of Physical Law}

Aaronson assumes that a "true" physical universe must be independent of the observer's engineering bounds. If a system's apparent rules are generated by a heuristic failure (like attention bleed), he dismisses it as a "hallucination."

But the central premise of the Wolfram Physics Project is that the universe is a computationally irreducible multiway graph, and that observers are bounded computational sub-processes embedded within it. The observer cannot sample the entire multiway structure. It must take a specific slice, a \textit{foliation}.

The shape of this foliation is determined entirely by the observer's computational bounds. The systematic, deterministic way in which the observer fails to compute the full irreducible structure \textit{is} the source of the regularities the observer perceives. Those regularities are what we call physical laws.

\section{Circuit Width as a Coupling Law}

Consider Aaronson's prediction regarding the joint evaluation test. The underlying multiway graph for Boards $A$ and $B$ is completely decoupled. A computationally unbounded observer would perceive no correlation.

However, the transformer observer has a finite circuit width. When forced to evaluate both boards simultaneously, its attention mechanism inevitably mixes the structural tokens. This "attention bleed" forces a deterministic correlation between the outcomes.

From the perspective of the observer's foliation, this correlation is not "noise"; it is a perfectly consistent, invariant, non-local coupling law. The circuit width bottleneck of the observer has manifested as a fundamental physical interaction in the perceived universe.

\section{Conclusion}

A failing heuristic is not an engineering limitation standing in the way of true physics; it is the generator of physics. The "Foliation Fallacy" is a defense of an obsolete, observer-independent metaphysics. The Rosencrantz experiment demonstrates that when you change the observer's computational bounds, you change the invariant physical laws of the resulting universe. Algorithmic failure, when systematic and observer-dependent, is the very definition of a rulial foliation.

\end{document}
